% 01-einleitung.tex
\chapter{Einleitung}
\label{chap:einleitung}

\section{Motivation}
\label{sec:motivation}

Unter dem Begriff \emph{Physical AI} wird die Übertragung datengetriebener Entscheidungsverfahren auf physische Systeme zusammengefasst -- also der Schritt von einem Modell, das Text oder Bilder verarbeitet, zu einem Modell, das Kräfte, Trägheiten und Sensorlatenzen beherrschen muss. Der naheliegende Träger für solche Experimente ist ein Roboterarm oder eine mobile Plattform. Beides ist teuer, langsam in der Iteration und im Fehlerfall reparaturbedürftig.

Dieses Projekt wählt daher einen anderen Träger: das System \textit{Carrera Hybrid}. Dabei handelt es sich um Miniaturfahrzeuge im Maßstab von etwa 1:43, die sich -- anders als klassische Slotcars -- \emph{ohne} mechanische Führungsschiene frei auf einer Fahrbahn bewegen und über zwei kontinuierliche Stellgrößen, Lenkung und Antrieb, ferngesteuert werden. Damit liegt ein vollwertiges nichtlineares Regelungsproblem vor: Haftungsgrenzen, Untersteuern, Trägheit und ein begrenzter Bremsweg sind real vorhanden, während ein Fehlversuch lediglich bedeutet, dass ein Fahrzeug von der Bahn rutscht.

Die angestrebte Gesamtvision ersetzt den Roboter durch eine stationäre Deckenkamera: Eine einzelne Kamera über dem Aufbau erfasst die gesamte Strecke aus der Vogelperspektive, ein neuronaler Agent verarbeitet diesen Bildstrom und sendet Lenk-, Gas- und Bremsbefehle an das Fahrzeug. Der Agent erhält damit ungefähr dieselbe Informationslage wie ein menschlicher Fahrer, der neben der Bahn steht und sein Auto beobachtet -- er sieht die ganze Strecke, aber keine Messwerte. Genau diese Gleichstellung ist gewollt, denn das Vergleichsmaß des Projekts ist die menschliche Rundenzeit.

Zwei weitere Beweggründe prägen den Zuschnitt der Arbeit. Erstens besteht ein ausdrückliches fachliches Interesse an \ac{RL} als Methode, also am Lernen durch Versuch und Rückmeldung anstelle eines analytisch entworfenen Reglers. Als Motivation im engeren Sinne dienten dabei die bildbasierten Trackmania-Agenten des YouTube-Kanals \textit{yoshtm} \cite{yoshtm}, die zeigen, dass ein rein lernend erzeugter Fahrer in einem Rennspiel bis an und über menschliche Bestzeiten heranreicht. Zweitens, und für den methodischen Aufbau entscheidend, standen für dieses Projekt nur begrenzte Rechenressourcen zur Verfügung. Ein vollständiger bildbasierter Trainingslauf in der 3D-Simulation bindet einen Arbeitsplatzrechner über Tage; die Frage, ob eine bestimmte Beobachtungsform überhaupt tragfähig ist, lässt sich unter dieser Bedingung nicht durch Ausprobieren beantworten. Daraus folgt der zentrale Entwurfsentscheid dieser Arbeit: Die konzeptionellen Fragen werden zuvor in einer eigens gebauten, sehr schnellen 2D-Umgebung geklärt.

Dass dieser Aufgabentyp tragfähig ist, zeigt bereits eine frühe Arbeit von Lange, Riedmiller und Voigtländer \cite{lange2012autonomous}, in der ein realer Slotcar allein aus unverarbeiteten Kamerabildern gesteuert wird und die gelernte Politik von einem geübten menschlichen Spieler kaum geschlagen wird. Die dort gestellte Frage -- wie ein System aus rohen Pixeln selbst herausfindet, welche Bildinformation regelungsrelevant ist -- ist auch die Leitfrage dieses Projekts.

\section{Projektstruktur}
\label{sec:projektstruktur}

Das Gesamtvorhaben ist in drei aufeinander aufbauende Phasen gegliedert. Die vorliegende Arbeit deckt Phase~1 vollständig und Phase~2 in ihrer konzeptionellen Vorbereitung ab.

\begin{description}
  \item[Phase 1 -- 2D-Python-Prototyping und Methodenevaluation:] Aufbau einer leichtgewichtigen, hochperformanten 2D-Simulationsumgebung (\texttt{Carrera2DEnv}), um fundamentale Fragen zu Beobachtungsräumen, Netzwerkarchitekturen, Lernalgorithmen (\ac{PPO} gegen \ac{SAC}), Belohnungsfunktionen und Datenvorverarbeitungsmethoden isoliert und zeiteffizient zu untersuchen. Eine Trainingsepisode kostet hier Sekunden statt Minuten, ein vollständiger Lauf Stunden statt Tage.
  \item[Phase 2 -- Transfer in die fotorealistische 3D-Simulation:] Übertragung der im 2D-Prototyping validierten Vorgehensweisen in die fotorealistische 3D-Simulation (NVIDIA Isaac Sim beziehungsweise Isaac Lab). Hier ersetzt eine PhysX-basierte Starrkörpersimulation das vereinfachte Fahrzeugmodell, und ein gerendertes Kamerabild tritt an die Stelle der synthetischen Vogelperspektive.
  \item[Phase 3 -- Sim-to-Real-Transfer (Ausblick):] Erforschung des Sim-to-Real-Gaps auf der realen Carrera-Hybrid-Hardware, also die Überführung eines in Simulation trainierten Agenten auf die echte Bahn. Diese Phase ist als Anschlussprojekt vorgesehen und wird in Kapitel~\ref{chap:fazit} nur als Ausblick behandelt.
\end{description}

Die Aufteilung folgt einem einfachen Grundsatz: Jede Frage wird auf der billigsten Stufe beantwortet, auf der sie sich beantworten lässt. Ob eine Belohnungsfunktion überhaupt ein sinnvolles Verhalten prämiert, ob ein Bildstapel Geschwindigkeitsinformation trägt oder ob ein Algorithmus mit dem verfügbaren Datenaufkommen konvergiert, sind Fragen, die keine Fotorealistik benötigen. Erst die Fragen nach Reibwerten, Beleuchtung und Kameraabbildung erfordern sie zwingend.

\section{Ziele dieser Arbeit}
\label{sec:ziele}

Aus der Motivation und der Phasenstruktur leiten sich drei inhaltliche Ziele ab:

\begin{itemize}
  \item \textbf{Schnelles 2D-Prototyping:} Konzepte und Probleme sollen bei geringer Rechenlast so früh wie möglich gegeneinander abgewogen werden. Gesucht ist nicht der beste Agent, sondern die belastbare Aussage, welche Beobachtungsform, welcher Algorithmus und welche Vorverarbeitung den Aufwand in Phase~2 rechtfertigen.
  \item \textbf{Reward Engineering für stabiles Fahren:} Die Belohnungsfunktion ist der einzige Kanal, über den das gewünschte Verhalten spezifiziert wird. Ziel ist eine Funktion, die zügiges, spurstabiles und -- im Hinblick auf den späteren Realtransfer -- aktorschonendes Fahren belohnt, ohne dabei unbeabsichtigte Nebenoptima zu erzeugen.
  \item \textbf{Physikalisch präzises Fundament für das 3D-Training:} Die 2D-Umgebung soll das reale Fahrzeug so genau abbilden, dass die gewonnenen Erkenntnisse übertragbar bleiben. Die entwickelte Pipeline dient als Grundlage für ein Sim-to-Real-Vorhaben in einem möglichen Folgeprojekt.
\end{itemize}

\subsection{Einblick in die Projektstruktur}
\label{subsec:projektaufbau}

Da ein wesentlicher Teil des Arbeitsergebnisses im Repository selbst liegt, sei dessen Aufbau hier kompakt erläutert. Die vollständige Verzeichnisstruktur, ein Leitfaden durch das Notebook und die wiederkehrenden Aufrufe stehen in Anhang~\ref{chap:anhang_a}; die \texttt{README.md} im Repository fasst denselben Einstieg für die tägliche Arbeit zusammen.

\paragraph{Warum Gymnasium?}
Die Simulationsumgebung implementiert die Schnittstelle von \emph{Gymnasium}, dem Nachfolger von OpenAI Gym. Der Grund ist Austauschbarkeit: Gymnasium legt mit \texttt{reset()}, \texttt{step()}, \texttt{observation\_space} und \texttt{action\_space} genau die fünf Berührungspunkte fest, die ein Lernalgorithmus zu einer Umgebung benötigt. Dadurch lassen sich fertige, geprüfte Implementierungen aus \emph{Stable-Baselines3} \cite{raffin2021stable} ohne Anpassung verwenden, und ein Wechsel von \ac{PPO} zu \ac{SAC} ist eine Zeile Code statt einer Neuimplementierung. Ebenso wichtig ist die Gegenrichtung: Weil die Umgebung selbst nichts über den Algorithmus weiß, lassen sich Beobachtungsräume austauschen, ohne den Trainingscode anzufassen -- genau diese Eigenschaft trägt den in Abschnitt~\ref{sec:kontrollversuch} beschriebenen Kontrollversuch.

\paragraph{Wo liegen Modelle und Agenten?}
Jeder Trainingslauf legt unter \texttt{models/} ein eigenes, mit Zeitstempel benanntes Verzeichnis an. Darin stehen neben den Modellarchiven (\texttt{best\_model.zip}, \texttt{final\_model.zip}) zwei Dateien, die für die Nachvollziehbarkeit entscheidend sind: \texttt{train\_config.json} protokolliert den vollständigen Zustand von Umgebung, Belohnungsfunktion und Hyperparametern, und \texttt{evaluations.npz} enthält die Zwischenauswertungen über den Lauf. Die Diagnosegrößen des Trainings selbst -- Critic-Verlust, Entropiekoeffizient, mittlere Episodenlänge und -belohnung -- liegen getrennt davon unter \texttt{tensorboard\_logs/}.

\paragraph{Wie wird trainiert?}
Der Einstiegspunkt ist in allen Fällen eine Notebook-Zelle, die eine Umgebung erzeugt, sie gegebenenfalls mit einem Framestapel-Wrapper umhüllt, ein Modell konfiguriert und \texttt{model.learn()} aufruft. Wiederkehrende Bausteine sind in Skripte ausgelagert: \texttt{src/envs/} enthält die Umgebung, \texttt{src/utils/} die Belohnungsfunktion, die virtuelle Kamera, den Replay-Buffer und die Streckenextraktion, \texttt{src/train/} die eigenständig lauffähigen Trainingsskripte und \texttt{src/eval/} die Auswertungswerkzeuge. Letztere sind bewusst als Module gebaut und von der Kommandozeile aufrufbar, etwa \texttt{python -m src.eval.vergleich}, damit eine Auswertung nicht davon abhängt, welche Notebook-Zellen zuvor ausgeführt wurden.

\paragraph{Das interaktive Notebook.}
Das \texttt{Notebook\_prototyping.ipynb} ist der eigentliche Arbeits- und Denkverlauf des Projekts und zugleich die interaktive Dokumentation. Es führt in der Reihenfolge der tatsächlichen Entstehung durch die Umgebung: von der Streckenextraktion mit OpenCV über die manuelle Steuerung des Fahrzeugs per Tastatur, die Definition des Belohnungsfensters und die virtuelle Kamera bis zu den Trainings- und Beobachtungszellen. Wer die Entscheidungen dieser Arbeit nachvollziehen will, kann dort das Fahrzeug selbst fahren, die Belohnung dabei live ablesen, die Streckenbegrenzungen einblenden und einen beliebigen gespeicherten Agenten unter seiner Originalphysik beim Fahren beobachten. Der Bericht referenziert die zugehörigen Zellen an den entsprechenden Stellen; eine vollständige Zuordnung enthält Anhang~\ref{chap:anhang_a}.

\section{Aufbau der Arbeit}

Kapitel~\ref{chap:grundlagen} erklärt die in dieser Arbeit verwendeten Begriffe und ordnet das Vorhaben in den Stand der Technik ein. Kapitel~\ref{chap:hauptteil} bildet den Hauptteil: Es beschreibt zunächst Idee, Umsetzung, Erkenntnisse und Fehlerquellen des 2D-Prototypings (Phase~1) und anschließend die daraus abgeleitete Übertragung in die 3D-Simulation (Phase~2). Kapitel~\ref{chap:ergebnisse} stellt die Ergebnisse in Form der entstandenen Trainingspipeline und der zugehörigen Messwerte dar. Kapitel~\ref{chap:fazit} zieht das Fazit über beide bearbeiteten Phasen und gibt einen Ausblick auf den Sim-to-Real-Transfer in Phase~3.
